\chapter{Datos}
\label{cap:datos}

Este capítulo describe de dónde vienen los datos, cómo se consolidan y qué contiene el conjunto
resultante. La descripción variable por variable está en el Apéndice~\ref{ap:variables}.

\section{Fuentes}
\label{sec:datos_fuentes}

\paragraph{Arribos pesqueros.} La variable objetivo proviene de dos fuentes complementarias.
La primera es el registro histórico de COBI (2016-2021), obtenido por convenio, con
granularidad de aviso de arribo. La segunda son los \emph{avisos de cosecha} de CONAPESCA
(``AVISOS MAYORES/MENORES DE COSECHA'', 2022-2026), publicados como datos abiertos
\citep{conapesca_datos}. La segunda fuente fue la
que permitió duplicar el número de temporadas disponibles.

\paragraph{Oceanografía.} La temperatura superficial del mar (SST) proviene del producto
reprocesado L4 de Copernicus Marine Service (OSTIA,
\texttt{METOFFICE-GLO-SST-L4-REP-OBS-SST}), con cobertura diaria global desde 1981. OSTIA es un
análisis que combina observaciones satelitales e in situ en una malla libre de nubes
\citep{donlon2012}, en la configuración operativa descrita por \citet{good2020}. Las
variables de color del océano ---clorofila (CHL), coeficiente de atenuación difusa (KD490),
materia particulada en suspensión (SPM), profundidad de disco de Secchi (ZSD), retrodispersión
de partículas (BBP) y materia orgánica coloreada disuelta (CDM)--- provienen del acervo
GlobColour \citep{globcolour2020} redistribuido por Copernicus, a 4~km de resolución.

\paragraph{Nota sobre el acceso a GlobColour.} La primera iteración de este trabajo descargaba
GlobColour por FTP. Esa vía entrega archivos \emph{globales sin recorte} (18~MB por variable y
por día a 4~km, del orden de cientos de gigabytes para el periodo y las variables de interés),
lo que la hace inviable. Se optó por la misma información redistribuida por Copernicus, que
permite recortar del lado del servidor. Las variables atmosféricas del acervo original (POC,
PIC, PAR, aerosoles, nubosidad) no están disponibles por esa vía y se omiten.

\section{Consolidación}
\label{sec:datos_etl}

Se reescribió el proceso ETL de forma modular y reproducible. Los arribos pesqueros se ingieren con granularidad de una fila por (fecha, especie, unidad económica) uniendo dos fuentes complementarias: el export histórico de COBI (2016-2021) y los avisos de cosecha (``AVISOS MAYORES/MENORES DE COSECHA'') de CONAPESCA (2022-2026). La unión se resuelve por prioridad temporal ---COBI hasta el 31-diciembre-2021 y CONAPESCA a partir de 2022, cada archivo acotado a su propio año--- de modo que no hay doble conteo en la costura de 2021. El lector de CSV autodetecta el encabezado (el preámbulo de los avisos varía entre 2 y 4 líneas) y la limpieza acepta rangos de fecha explícitos. El conjunto consolidado cubre así 2017-2026, casi duplicando las temporadas disponibles de langosta en San Quintín (de $\sim$5 a $\sim$9). Como fuente de temperatura superficial del mar (SST) se adoptó el producto reprocesado L4 de Copernicus Marine Service (OSTIA, \texttt{METOFFICE-GLO-SST-L4-REP-OBS-SST}), con cobertura diaria global desde 1981; se descargó la serie 1982-2025 y se promedió sobre el bounding box costero de cada unidad económica. Esta fuente reemplaza al cliente \texttt{motuclient} —descontinuado en 2024— por el nuevo SDK \texttt{copernicusmarine}, y provee la profundidad histórica necesaria para construir una climatología de referencia de 30 años.


\paragraph{Consistencia del empalme.} Como las dos fuentes de arribos se unen en la frontera de
2021, se verificó que la costura no introdujera un salto de escala. No hay fechas duplicadas;
COBI termina el 31-diciembre-2021 ---truncando la temporada 2021-2022 a unas 31~t--- y la unión
recupera enero y febrero de 2022, completando esa temporada a 39.5~t. La transición es continua:
diciembre de 2021 según COBI (6.3~t) es comparable con enero de 2022 según CONAPESCA (6.5~t).
La limitación honesta es que \emph{no existe un periodo de solape} entre ambas fuentes que
permita validar la escala de forma exhaustiva; el traspaso diciembre$\to$enero es la evidencia
disponible y es consistente.

\paragraph{Esquema.} El conjunto consolidado tiene una fila por
$(\text{fecha}, \text{especie}, \text{\ue{}})$ sobre el rango de fechas completo ---no solo los
días con registro---, con la variable objetivo, las covariables oceanográficas, el indicador de
temporada, la temporada a la que pertenece el día, banderas de imputación y metadatos de
trazabilidad del proceso. Se almacena en formato Parquet particionado por especie y año.

\section{Tratamiento de la variable objetivo faltante}
\label{sec:datos_faltantes}

Como el conjunto cubre todos los días y no solo los días con aviso, hay que decidir qué valor
toma la captura en los huecos. La semántica importa: un $0$ afirma ``no se pescó'' y un valor
faltante afirma ``no sabemos''. La decisión adoptada distingue tres casos:

\begin{enumerate}
  \item \textbf{Fuera de temporada} sin registro $\Rightarrow y = 0$. No es una imputación sino
  la consecuencia lógica de la veda.
  \item \textbf{Dentro de temporada} sin registro $\Rightarrow y$ faltante. Puede ser un día sin
  pesca (clima, festivo) o un hueco de captura de datos, y no hay forma de distinguirlos, así
  que no se inventa un valor.
  \item Cualquier imputación explícita queda marcada con una bandera, de modo que sea
  rastreable. No se hace por omisión.
\end{enumerate}

Cada modelo decide cómo tratar el faltante en su propia capa de preparación ---en los
experimentos de este trabajo se rellena con cero dentro de temporada, bajo el supuesto de que un
día en temporada sin registro es un día sin captura--- pero el conjunto de datos preserva la
distinción. La decisión está documentada como ADR en el repositorio.

\section{Zona de cada \ue{}}
\label{sec:datos_zona}

La oceanografía se agrega espacialmente sobre la zona costera de cada \ue{}. Lo correcto sería
recortar con el polígono TURF exacto de cada cooperativa, pero esos polígonos no son de acceso
público y no fue posible obtenerlos. Se usa en su lugar un rectángulo costero (\emph{bounding
box}) por clúster de cooperativas. A la resolución de los productos satelitales empleados
---0.25\textdegree{} para la SST y 4~km para el color del océano--- el rectángulo es una
aproximación razonable, y es además lo que hacía la primera iteración del trabajo; el polígono
exacto queda como mejora deseable y no bloqueante. La Figura~\ref{fig:mapa_ues} muestra las
zonas sobre el campo medio de SST, y la Figura~\ref{fig:zona_pesca} el detalle de la zona de
San Quintín.

\begin{figure}[htbp]
\centering
\includegraphics[width=0.78\linewidth]{mapa_ues.png}
\caption{Zonas costeras asignadas a los clústeres de \UEs{} del conjunto consolidado, sobre la
SST media de 2025. La península de Baja California aparece en blanco (sin dato de SST sobre
tierra) y el agua cálida al oriente es el Golfo de California. El gradiente térmico a lo largo
del rango es explícito: de $\sim$17\textdegree{}C en San Quintín ($\sim$30.4\textdegree{}N) a
$\sim$22\textdegree{}C en Bahía Magdalena ($\sim$24.5\textdegree{}N).}
\label{fig:mapa_ues}
\end{figure}

\section{Índice de olas de calor marinas}
\label{sec:datos_mhw}

Se implementó el índice de MHW siguiendo la metodología de Hobday et al. \cite{hobday2016}: para cada día del año se calcula una climatología y su percentil 90 sobre el periodo base 1982-2011; un evento MHW es un tramo de al menos 5 días consecutivos con SST por encima de ese umbral (fusionando eventos separados por huecos de $\leq 2$ días), y se clasifica en categorías I-IV según cuántos múltiplos de la diferencia (umbral $-$ media) alcanza la anomalía. El dataset consolidado incorpora tres columnas nuevas: \texttt{sst\_anomaly} (siempre), \texttt{mhw\_category} (0-4) y \texttt{mhw\_intensity} (definida durante eventos).

Como validación, el índice reproduce con nitidez los eventos conocidos del Pacífico nororiental: el llamado ``Blob'' de 2014-2016 (con 231, 293 y 102 días de MHW respectivamente en la zona de San Quintín) y el régimen cálido de 2019-2021 (Figura \ref{fig:mhw_timeline}). Esta validación es relevante porque conecta directamente con el mecanismo documentado por \cite{villasenor2024}.

\begin{figure}[htbp]
\centering
\includegraphics[width=\linewidth]{mhw_timeline.png}
\caption{Serie diaria de SST en el bbox de San Quintín (negro) contra su climatología (azul continuo) y el umbral del percentil 90 (azul discontinuo), con baseline 1982-2011 (Hobday et al. 2016). Las áreas sombreadas marcan los eventos de ola de calor marina. Destacan el ``Blob'' de 2014-2016 y el régimen cálido de 2019-2021 que precede al colapso de la captura de langosta de la temporada 2021-2022.}
\label{fig:mhw_timeline}
\end{figure}


\section{Construcción de variables}
\label{sec:datos_features}

A partir del conjunto consolidado se construye la matriz de variables explicativas de la
ecuación~\eqref{eq:problema}, con tres bloques:

\begin{itemize}
  \item \textbf{Calendario} (determinista en $t$): día del año codificado como $(\sin, \cos)$ e
  indicador de temporada.
  \item \textbf{Rezagos de la propia serie}: la captura de la misma fecha de uno y dos años
  atrás, que captura el ciclo anual.
  \item \textbf{Oceanografía desplazada}: cada covariable oceanográfica en $t-\horizonte$, más
  medias móviles de 90 y 365 días que terminan en $t-\horizonte$ ---es decir, el estado del mar
  previo a la temporada y el calor acumulado del año anterior. El desplazamiento es el mecanismo
  ecológico que se quiere modelar y, a la vez, lo que garantiza que no haya fuga de información
  (Figura~\ref{fig:desplazamiento}).
\end{itemize}

Con las seis variables de color del océano el conjunto crece a unas 35 variables, un número
grande frente a las pocas temporadas disponibles; el Capítulo~\ref{cap:resultados} muestra que
ese desbalance importa y cómo se corrige con selección de variables.

\section{El conjunto consolidado}
\label{sec:datos_resumen}

El resultado cubre \textbf{2017-2026} y contiene \textbf{44 series} (especie $\times$ \ue{}) con
al menos un día de captura registrado, que suman \textbf{10{,}465 toneladas} desembarcadas: 21
series de langosta roja (7{,}569~t), 14 de abulón azul, 6 de erizo rojo y 3 de otras especies de
abulón. La mediana de días con captura por serie es de \textbf{70}, lo que ilustra el problema
de fondo: aun con casi una década de cobertura, cada serie individual aporta pocas
observaciones. La Tabla~\ref{tab:series_resumen} desglosa cada serie.

% Generado por experiments/thesis_assets.py — no editar a mano.
\begin{longtable}{llrrrrc}
\caption{Series (especie $\times$ unidad económica) del conjunto consolidado 2017-2026. `Temporadas' cuenta las que tienen al menos un día de captura; la cobertura es la fracción de días con el valor oceanográfico disponible.}\label{tab:series_resumen}\\
\toprule
Especie & Zona & Temporadas & Días con captura & Toneladas & Cobertura SST & Cobertura CHL \\
\midrule
\endfirsthead
\toprule
Especie & Zona & Temporadas & Días con captura & Toneladas & Cobertura SST & Cobertura CHL \\
\midrule
\endhead
\bottomrule
\endfoot
Abulón negro & San Quintín (LITORAL DE BAJA CALIFO) & 5 & 51 & 2.5 & 100\% & 100\% \\
Abulón negro & El Rosario (SCPP ENSENADA SCL) & 5 & 20 & 2.3 & 100\% & 100\% \\
Abulón azul & Vizcaíno (SCPP PESCADORES NACION) & 4 & 33 & 175.1 & 100\% & 100\% \\
Abulón azul & El Rosario (SCPP ENSENADA SCL) & 9 & 68 & 59.7 & 100\% & 100\% \\
Abulón azul & Pacífico BCS (SCPP LA PURISIMA SC DE) & 4 & 28 & 24.8 & 100\% & 100\% \\
Abulón azul & Bahía Magdalena (SCPP PUERTO CHALE SCL) & 4 & 15 & 16.2 & 100\% & 100\% \\
Abulón azul & Vizcaíno (SCPP BAHIA TORTUGAS SC) & 1 & 8 & 9.3 & 100\% & 100\% \\
Abulón azul & Punta Canoas (SPR PUNTA CANOAS S DE ) & 9 & 31 & 7.8 & 100\% & 100\% \\
Abulón azul & San Quintín (LITORAL DE BAJA CALIFO) & 5 & 62 & 5.5 & 100\% & 100\% \\
Abulón azul & Pacífico BCS (SC PROGRESO DE PRODUCC) & 4 & 17 & 5.1 & 100\% & 100\% \\
Abulón azul & Bahía Magdalena (SCPP BAHIA MAGDALENA S) & 1 & 6 & 4.2 & 100\% & 100\% \\
Abulón azul & Pacífico BCS (SCPP CALIFORNIA DE SAN) & 1 & 5 & 2.6 & 100\% & 100\% \\
Abulón azul & Vizcaíno (SCPP BUZOS Y PESCADORE) & 3 & 4 & 2.6 & 100\% & 100\% \\
Abulón azul & Vizcaíno (SCPP EMANCIPACION SC D) & 4 & 12 & 2.1 & 100\% & 100\% \\
Abulón azul & Pacífico BCS (SCPP PUNTA ABREOJOS SC) & 1 & 2 & 0.3 & 100\% & 100\% \\
Abulón azul & El Rosario (SPR ASOCIACION PESQUER) & 1 & 1 & 0.0 & 100\% & 100\% \\
Abulón rojo & San Quintín (LITORAL DE BAJA CALIFO) & 5 & 42 & 2.9 & 100\% & 100\% \\
Langosta roja & El Rosario (SCPP ENSENADA SCL) & 10 & 688 & 1,181.2 & 100\% & 100\% \\
Langosta roja & Vizcaíno (SCPP PESCADORES NACION) & 10 & 307 & 904.2 & 100\% & 100\% \\
Langosta roja & Vizcaíno (SCPP BAHIA TORTUGAS SC) & 5 & 193 & 767.1 & 100\% & 100\% \\
Langosta roja & San Quintín (LITORAL DE BAJA CALIFO) & 9 & 1110 & 689.7 & 100\% & 100\% \\
Langosta roja & Vizcaíno (SCPP EMANCIPACION SC D) & 5 & 182 & 615.8 & 100\% & 100\% \\
Langosta roja & Pacífico BCS (SCPP LA PURISIMA SC DE) & 5 & 78 & 532.9 & 100\% & 100\% \\
Langosta roja & Vizcaíno (SCPP LEYES DE REFORMA ) & 6 & 148 & 493.8 & 100\% & 100\% \\
Langosta roja & Pacífico BCS (SC PROGRESO DE PRODUCC) & 6 & 133 & 366.5 & 100\% & 100\% \\
Langosta roja & Vizcaíno (SCPP BUZOS Y PESCADORE) & 5 & 143 & 362.1 & 100\% & 100\% \\
Langosta roja & El Rosario (PESCADORES EL CHUTE SP) & 10 & 307 & 282.4 & 100\% & 100\% \\
Langosta roja & El Rosario (ASOCIACION PESQUERA MO) & 10 & 269 & 267.7 & 100\% & 100\% \\
Langosta roja & Pacífico BCS (SCPP PUNTA ABREOJOS SC) & 5 & 109 & 260.2 & 100\% & 100\% \\
Langosta roja & El Rosario (SPR ASOCIACION PESQUER) & 10 & 392 & 232.4 & 100\% & 100\% \\
Langosta roja & El Rosario (ISLA SAN GERONIMO SPR ) & 10 & 316 & 162.3 & 100\% & 100\% \\
Langosta roja & Pacífico BCS (SCPP CALIFORNIA DE SAN) & 5 & 67 & 152.0 & 100\% & 100\% \\
Langosta roja & Punta Canoas (SPR PUNTA CANOAS S DE ) & 10 & 208 & 82.6 & 100\% & 100\% \\
Langosta roja & San Quintín (EL PABELLON DE SAN QUI) & 8 & 193 & 80.0 & 100\% & 100\% \\
Langosta roja & San Quintín (ROCAS DE SAN MARTIN SP) & 10 & 217 & 58.3 & 100\% & 100\% \\
Langosta roja & Bahía Magdalena (SCPP PUERTO CHALE SCL) & 5 & 44 & 34.2 & 100\% & 100\% \\
Langosta roja & Bahía Magdalena (SCPP BAHIA MAGDALENA S) & 6 & 72 & 25.3 & 100\% & 100\% \\
Langosta roja & Bahía Magdalena (SCPP PUERTO SAN CARLOS) & 5 & 89 & 18.8 & 100\% & 100\% \\
Erizo rojo & El Rosario (SCPP ENSENADA SCL) & 10 & 490 & 1,483.3 & 100\% & 100\% \\
Erizo rojo & El Rosario (SPR ASOCIACION PESQUER) & 10 & 487 & 644.6 & 100\% & 100\% \\
Erizo rojo & El Rosario (ASOCIACION PESQUERA MO) & 10 & 406 & 406.5 & 100\% & 100\% \\
Erizo rojo & El Rosario (PESCADORES EL CHUTE SP) & 4 & 21 & 19.2 & 100\% & 100\% \\
Erizo rojo & El Rosario (ISLA SAN GERONIMO SPR ) & 2 & 9 & 10.4 & 100\% & 100\% \\
Erizo rojo & Punta Canoas (SPR PUNTA CANOAS S DE ) & 9 & 23 & 8.4 & 100\% & 100\% \\
\end{longtable}


\section{Origen de la información de las figuras heredadas}

La primera iteración de este trabajo construyó su proceso de datos sobre GlobColour por FTP y
sobre un intercambio directo de archivos con COBI. Ese proceso se muestra en la
Figura~\ref{fig:etl_globcolour} y se conserva porque documenta el punto de partida; el proceso
vigente es el descrito en la Sección~\ref{sec:datos_etl}.

\begin{figure}[htbp]
    \centering
    \includegraphics[width=0.95\linewidth]{Imagen1.png}
    \caption{Proceso de datos de la primera iteración del trabajo: ingesta de GlobColour por FTP
    y de los archivos de COBI, transformación de \texttt{.nc} a tablas y consolidación. El
    proceso vigente, descrito en la Sección~\ref{sec:datos_etl}, sustituye la ingesta por FTP
    por el SDK de Copernicus con recorte del lado del servidor.}
    \label{fig:etl_globcolour}
\end{figure}

